% ════════════════════════════════════════════════════════════════
\subsection{Egzamin zerowy}
% ════════════════════════════════════════════════════════════════

\begin{zadanie}[Zadanie 1 --- dwa najtańsze symbole pod wspólnym rodzicem (15 p.)]
	Udowodnij, że w~kodowaniu Huffmana istnieje optymalne drzewo kodowe,
	w~którym dwa najtańsze (najrzadziej występujące) symbole mają wspólnego
	rodzica.
\end{zadanie}

To jest dokładnie Lemat~\ref*{N-lemma:huffman:1} z~wykładu; przepiszemy jego dowód poniżej

\begin{proof}
	Znajdziemy drzewo reprezentujące optymalny kod prefiksowy, w którym liście reprezentujące $x$ i $y$ mają wspólnego rodzica i mają największą głębokość w drzewie.

	Niech $T$ będzie dowolnym drzewem reprezentującym optymalny kod prefiksowy i niech $a$ i $b$ będą dwoma liśćmi o tym samym rodzicu i najmniejszej głębokości w $T$.
	Bez straty ogólności możemy przyjąć, że $f(x) \leq f(y)$ i $f(a) \leq f(b)$.
	$f(x)$ i $f(y)$ są najmniejszymi kluczami, więc $f(x) \leq f(a)$ i $f(y) \leq f(b)$.

	Niech $T'$ będzie drzewem powstałym z $T$ przez zamianę liści $x$ i $a$.
	\begin{flalign*}
		B(T) - B(T') & = \sum_{c \in C} f(c) \cdot d_T(c) - \sum_{c \in C} f(c) \cdot d_{T'}(c) =                                        &  & \\
		             & = \vphantom{\sum_{c \in C}} f(a) \cdot d_T(a) + f(x) \cdot d_T(x) - f(a) \cdot d_{T'}(a) - f(x) \cdot d_{T'}(x) = &  & \\
		             & = \vphantom{\sum_{c \in C}} f(a) \cdot d_T(a) + f(x) \cdot d_T(x) - f(a) \cdot d_T(x) - f(x) \cdot d_T(a) =       &  & \\
		             & = \vphantom{\sum_{c \in C}} f(a) (d_T(a) - d_T(x)) - f(x) (d_T(a) - d_T(x)) =                                     &  & \\
		             & = \vphantom{\sum_{c \in C}} \underbrace{(f(a) - f(x))}_{\geq 0} \underbrace{(d_T(a) - d_T(x))}_{\geq 0} \geq 0    &  &
	\end{flalign*}

	Zatem $B(T) \geq B(T')$, ale $B(T)$ najmniejsze możliwe, więc $B(T') = B(T)$ czyli $T'$ też reprezentuje optymalny kod prefiksowy. Analogicznie zamieniamy miejscami liście $y$ i $b$ i otrzymujemy drzewo $T''$ reprezentujące optymalny kod prefiksowy, które ma żądane własności.
\end{proof}

\clearpage
\begin{zadanie}[Zadanie 2 --- skojarzenie najliczniejsze a~ścieżki powiększające (15 p.)]
	Udowodnij, że skojarzenie w~grafie jest najliczniejsze wtedy i~tylko wtedy,
	gdy nie istnieje ścieżka powiększająca.
\end{zadanie}

Jest to dokładnie twierdzenie Berge'a (twierdzenie~\ref*{N-theorem:berge} z~notatek). Przepiszemy dowód poniżej.

\begin{proof}
	Niech $M$ będzie skojarzeniem w~$G$.
	\begin{itemize}
		\item[$\implies$]
		      Przypuśćmy, że istnieje ścieżka powiększająca $P$ względem $M$. \\
		      Wtedy $M' = M \div E(P)$ jest skojarzeniem i $|M'| > |M|$, więc $M$ nie jest skojarzeniem najliczniejszym.
		\item[$\impliedby$]
		      Przypuśćmy, że $M$ nie jest najliczniejszym skojarzeniem w $G$. \\
		      Niech $M'$ będzie skojarzeniem takim, że $|M'| > |M|$. \\
		      Niech $H$ będzie podgrafem $G$ indukowanym przez $M' \div M$, czyli $H = (V(G), M' \div M)$. \\
		      $M, M'$ to skojarzenia, więc stopień wierzchołka $d_H(v) \le 2$ dla każdego $v \in V(G)$.

		      Każda składowa $H$ jest ścieżką lub cyklem. \\
		      $M, M'$ to skojarzenia, więc nie ma cykli nieparzystych w $H$. \\
		      Każda składowa $H$ jest ścieżką lub cyklem parzystym. \\
		      Skoro $|M'| > |M|$, to istnieje składowa $H$, która jest ścieżką zawierającą więcej krawędzi z $M'$ niż z $M$, więc ta ścieżka zaczyna się i kończy krawędzią z $M'$. Ta ścieżka jest ścieżką powiększającą względem $M$. \qedhere
	\end{itemize}
\end{proof}

\clearpage
\begin{zadanie}[Zadanie 3 --- matroid jednego wierzchołka na składową (15 p.)]
	Niech $G$ będzie grafem, a~$\mathcal{A}_G$ rodziną podzbiorów $V(G)$ taką, że
	$S\in\mathcal{A}_G$ wtedy i~tylko wtedy, gdy każda składowa spójności $G$
	zawiera co najwyżej jeden wierzchołek z~$S$. Rozstrzygnij, czy para
	$(V(G),\mathcal{A}_G)$ zawsze jest matroidem.
\end{zadanie}

\paragraph{Tak --- jest to \emph{partition matroid}.}
\begin{dygresja}
	\uemph{Partition matroid} (matroid podziałowy) powstaje z~podziału zbioru bazowego $S$ na rozłączne
	klasy $S=C_1\cup\dots\cup C_p$ i~liczb $d_1,\dots,d_p\ge 0$ (pojemności): zbiór
	$A\subseteq S$ jest niezależny $\iff |A\cap C_i|\le d_i$ dla każdej klasy $i$.
	Innymi słowy --- z~każdej klasy wolno wziąć co najwyżej $d_i$ elementów,
	niezależnie od~pozostałych klas. Że to faktycznie matroid, sprawdzamy poniżej;
	tutaj wszystkie $d_i=1$, a~klasami są składowe spójności $G$.

	\emph{Suma prosta matroidów jednorodnych.} Każda klasa $C_i$ z~pojemnością
	$d_i$ to osobny \emph{uniform matroid} $U^{d_i}_{|C_i|}$ (niezależne są zbiory
	$\le d_i$ elementów), a~cały partition matroid jest ich \uemph{sumą prostą}.
	W~szczególności \emph{uniform matroid} to przypadek jednej klasy $C_1=S$
	z~pojemnością $r$, oznaczany $U^r_n$ ($n=|S|$).

	\emph{Baza, cykl, rząd.} \uemph{Baza} bierze z~każdego bloku dokładnie $d_i$
	elementów; \uemph{cykl} (minimalny zbiór zależny) to $d_i{+}1$ elementów
	z~jednego bloku; \uemph{rząd} matroidu to $\sum_i d_i$. (Dla naszego $G$:
	baza = jeden wierzchołek na~składową, rząd = liczba składowych.)

	\emph{Powiązania.} Klasa matroidów podziałowych jest zamknięta na~dualność,
	minory i~sumy prostych. Gdy wszystkie $d_i=1$, partition matroid jest zarazem
	\emph{transversal matroid} rodziny rozłącznych bloków; na~tym opiera się ujęcie
	skojarzeń w~grafie dwudzielnym $(U,V)$ jako \uemph{przecięcia dwóch matroidów}
	(po~jednym ,,brak wspólnego końca'' dla strony $U$ i~dla $V$).
\end{dygresja}
Niech $C_1,\dots,C_p$ będą składowymi spójności $G$; tworzą one podział
$V(G)$. Z~definicji $S\in\mathcal{A}_G\iff |S\cap C_i|\le 1$ dla każdego $i$.
Sprawdzamy aksjomaty:

\newlength\labelw
\settowidth\labelw{(wł.\ 1)\ }
\begin{enumerate}[
		label=(wł.\ \arabic*),
		align=left,
		labelwidth=\labelw,
		labelsep=0pt,
		leftmargin=\dimexpr\labelw-2em\relax,
		itemindent=2em
	]
	\item[(wł.\ \ref*{N-matroid:skończoność})] $V(G)$ jest skończony i~niepusty. $\checkmark$
	\item[(wł.\ \ref*{N-matroid:niepustość})] $\emptyset\in\mathcal{A}_G$, więc $\mathcal{A}_G\neq\emptyset$. $\checkmark$
	\item[(wł.\ \ref*{N-matroid:dziedziczność})] Jeśli $S\in\mathcal{A}_G$ i~$S'\subseteq S$, to
	      $|S'\cap C_i|\le|S\cap C_i|\le 1$, więc $S'\in\mathcal{A}_G$.
	\item[(wł.\ \ref*{N-matroid:wymiana})] Niech $S_1,S_2\in\mathcal{A}_G$, $|S_1|<|S_2|$. Każda
	      składowa zawiera $\le 1$ element każdego ze zbiorów, więc $|S_j|$ równa
	      się liczbie składowych ,,trafionych'' przez $S_j$. Z~$|S_2|>|S_1|$
	      wynika, że istnieje składowa $C_i$ trafiona przez $S_2$, ale nie przez
	      $S_1$; biorąc $v\in S_2\cap C_i$ mamy $S_1\cup\{v\}\in\mathcal{A}_G$.
\end{enumerate}
Zatem $(V(G),\mathcal{A}_G)$ jest matroidem. $\qed$

\clearpage
\begin{zadanie}[Zadanie 4 --- dokładne pokrycie wierzchołkowe w~czasie podwykładniczym (15 p.)]
	Zaprojektuj algorytm dokładny, który znajdzie liczność najmniejszego pokrycia
	wierzchołkowego grafu w~czasie $o(2^n)$, gdzie $n$ jest liczbą wierzchołków.
\end{zadanie}

\paragraph{Rozgałęzianie po~wierzchołku dużego stopnia.}
Korzystamy z~obserwacji: dla dowolnej krawędzi pokrycie $C$ musi zawierać jeden
z~jej końców, a~równoważnie --- dla dowolnego wierzchołka $v$ albo $v\in C$,
albo (jeśli $v\notin C$) wszyscy jego sąsiedzi $N(v)\subseteq C$.

\begin{alg}{Minimalne pokrycie wierzchołkowe przez rozgałęzianie}{alg:minvc}
	\Function{MinVC}{$G$}\Comment{zwraca liczność najmniejszego pokrycia}
	\If{$G$ nie ma krawędzi}\State \Return $0$
	\EndIf
	\If{maksymalny stopień w~$G$ wynosi $\le 2$}
	\State \Return minimalne pokrycie sumy ścieżek i~cykli
	       \Comment{DP po~każdej ścieżce/cyklu}
	       \Statex \hfill czas wielomianowy
	\EndIf
	\State wybierz wierzchołek $v$ o~stopniu $\ge 3$
	\State $a\gets 1+{}$\Call{MinVC}{$G-v$}\Comment{gałąź $v\in C$}
	\State $b\gets |N(v)|+{}$\Call{MinVC}{$G-v-N(v)$}\Comment{gałąź $v\notin C$, więc $N(v)\subseteq C$}
	\State \Return $\min(a,b)$
	\EndFunction
\end{alg}

\begin{dygresja}
\uemph{DP ({\itshape Dynamic Programming}) po~ścieżce/cyklu.} Przy maksymalnym stopniu $\le 2$ graf rozpada
	się na~rozłączne ścieżki i~cykle; każdą składową liczymy osobno i~sumujemy
	wyniki. Dla ścieżki $v_1,\dots,v_k$ trzymamy dwa stany:
	\begin{itemize}[nosep]
		\item $f[i][0]$ --- minimalne pokrycie wierzchołków $v_1,\dots,v_i$ przy
		      $v_i\notin C$;
		\item $f[i][1]$ --- to samo przy $v_i\in C$.
	\end{itemize}
	Każda krawędź $v_{i-1}v_i$ musi być pokryta, stąd przejścia
	\[
		f[i][0]=f[i-1][1],\qquad
		f[i][1]=1+\min\big(f[i-1][0],\,f[i-1][1]\big).
	\]
	Jeśli $v_i\notin C$, to poprzednik musi być w~$C$ (inaczej krawędź
	$v_{i-1}v_i$ niepokryta); jeśli $v_i\in C$, poprzednik dowolny. Baza
	to~$f[1][0]=0$, $f[1][1]=1$, a~wynik ścieżki to $\min\big(f[k][0],f[k][1]\big)$.

	Dla cyklu dochodzi jedna domykająca krawędź $v_k v_1$. Robimy ten sam
	rachunek co dla ścieżki $v_1,\dots,v_k$, ale uruchamiamy go dwukrotnie,
	ustalając z~góry stan $v_1$ --- różnica jest tylko w~\emph{bazie} (w~$v_1$)
	i~w~\emph{odczycie} (w~$v_k$):
	\begin{itemize}[nosep]
		\item \emph{wymuszamy $v_1\in C$:} baza $f[1][1]=1$, $f[1][0]=+\infty$;
		      krawędź $v_k v_1$ jest pokryta przez $v_1$, więc $v_k$ dowolny ---
		      odczyt $\min\big(f[k][0],f[k][1]\big)$;
		\item \emph{wymuszamy $v_1\notin C$:} baza $f[1][0]=0$, $f[1][1]=+\infty$;
		      teraz $v_k v_1$ wymusza $v_k\in C$ --- odczyt $f[k][1]$.
	\end{itemize}
	(Wartość $+\infty$ zakazuje stanu sprzecznego z~wymuszeniem.) Każda
	składowa w~czasie liniowym, więc cały ten przypadek jest wielomianowy.
\end{dygresja}

\paragraph{Poprawność.}
Przy stopniu $\ge 3$ każde pokrycie wpada do~jednej z~dwóch gałęzi (zawiera $v$
lub całe $N(v)$), a~obie gałęzie zwracają poprawne minimum dla swojej
podinstancji; bierzemy mniejsze. Przypadek maksymalnego stopnia $\le 2$
rozwiązujemy wprost (graf jest sumą ścieżek i~cykli).

\paragraph{Złożoność.}
Gałąź pierwsza usuwa $1$ wierzchołek, druga co najmniej $1+3=4$ (sam $v$ oraz
$\ge 3$ jego sąsiadów). Składnik $\mathrm{poly}(n)$ to praca poza wywołaniami
rekurencyjnymi: znalezienie wierzchołka o~stopniu $\ge 3$ (lub stwierdzenie, że
maksymalny stopień $\le 2$), zbudowanie podgrafów $G-v$ oraz $G-v-N(v)$ i~DP
po~ścieżkach/cyklach w~przypadku bazowym --- każda z~tych operacji jest
wielomianowa względem liczby wierzchołków i~krawędzi. Stąd
\[
	T(n)\le T(n-1)+T(n-4)+\mathrm{poly}(n).
\]
Równanie charakterystyczne $x^4=x^3+1$ ma pierwiastek dodatni
$\alpha\approx 1{,}3803$, więc $T(n)=\BO(\alpha^n)=\BO(1{,}39^n)=o(2^n)$. $\qed$

\clearpage
\begin{zadanie}[Zadanie 5 --- 3-aproksymacja problemu 3DC (15 p.)]
	Rozważmy problem --- nazwijmy go roboczo \emph{3DC} --- w~którym mamy daną
	rodzinę $\mathcal{R}$ 3-elementowych podzbiorów pewnego zbioru $X$
	i~poszukujemy jak najmniej licznego zbioru $S\subseteq X$ takiego, że każdy
	zbiór z~$\mathcal{R}$ zawiera przynajmniej jeden element z~$S$. Zaprojektuj
	wielomianowy algorytm 3-aproksymacyjny dla problemu 3DC (i~udowodnij, że stała
	aproksymacji istotnie jest nie większa niż~3). Możesz zainspirować się
	algorytmem 2-aproksymacyjnym dla problemu najmniejszego pokrycia
	wierzchołkowego.
\end{zadanie}

\paragraph{Algorytm (jak 2-aproksymacja dla pokrycia wierzchołkowego).}
\begin{alg}{3-aproksymacja problemu 3DC przez maksymalne dopasowanie}{alg:3dc}
	\Function{Approx3DC}{$X,\mathcal{R}$}
	\State $S\gets\emptyset$
	\While{istnieje $R\in\mathcal{R}$ rozłączny z~$S$}
	\State wybierz taki $R$
	\State $S\gets S\cup R$\Comment{dodaj \uemph{wszystkie} trzy elementy}
	\EndWhile
	\State \Return $S$
	\EndFunction
\end{alg}
Po~zakończeniu każdy zbiór z~$\mathcal{R}$ jest przecięty przez $S$, więc $S$ jest
poprawnym rozwiązaniem.

\paragraph{Współczynnik aproksymacji.}
Niech $R_1,\dots,R_t$ będą zbiorami wybranymi w~kroku~2. Są one
\uemph{parami rozłączne}: gdy wybieramy $R_i$, jest on rozłączny z~bieżącym $S$,
a~ten zawiera wszystkie elementy wcześniej wybranych $R_1,\dots,R_{i-1}$; zatem
$R_i$ jest rozłączny z~każdym $R_j$, $j<i$. Każde rozwiązanie optymalne $S^\ast$
musi przecinać każdy $R_i$, a~skoro zbiory $R_i$ są rozłączne, potrzebuje na to
co najmniej $t$ różnych elementów: $|S^\ast|\ge t$. Nasze rozwiązanie ma
$|S|=3t\le 3|S^\ast|$. Algorytm jest więc $3$-aproksymacyjny.

\paragraph{Złożoność.}
Każdy obrót pętli dokłada do~$S$ trzy nowe elementy (wybrany $R$ jest rozłączny
z~$S$), więc pętla wykonuje się co najwyżej $\lfloor |X|/3\rfloor$ razy. W~każdym
obrocie przeglądamy rodzinę $\mathcal{R}$ w~poszukiwaniu zbioru rozłącznego
z~$S$, sprawdzając dla każdego $R\in\mathcal{R}$ jego trzy elementy --- to
$O(|\mathcal{R}|)$ pracy. Łącznie $O(|X|\cdot|\mathcal{R}|)$, czyli czas
wielomianowy. $\qed$
